\chapter*{Executive summary}

\noindent
\begin{longtable}{|p{\dimexpr \linewidth-2\tabcolsep-2\arrayrulewidth}|}
\hline%------------------------------------------------------------
\sumheading  Title of Project \\
\hline%------------------------------------------------------------
 The Design of an Automated, Solar-Powered Greenhouse for Optimal Plant Growth \\[1ex]

\hline%------------------------------------------------------------
\sumheading  Objectives \\
\hline%------------------------------------------------------------
 Design and build a small-scale greenhouse capable of automating all the necessary elements to sustain amicable living conditions for various plants while enabling the user with wireless monitoring capabilities. \\[1ex]

\hline%------------------------------------------------------------
\sumheading  What is current practice and what are its limitations? \\
\hline%------------------------------------------------------------
 Current domestically applicable greenhouses rely heavily on active monitoring and care, this requires significant labour and time and can still not achieve ideal living conditions for some plants.  \\[1ex]

\hline%------------------------------------------------------------
\sumheading  What is new in this project? \\
\hline%------------------------------------------------------------
This project introduces a system capable of wirelessly monitoring multiple parameters within the greenhouse, including temperature, humidity, soil moisture, air composition, and airflow, and enables effective control of these parameters within user-defined ranges. The integration of solar power helps to offset the system’s energy requirements, reinforcing its sustainability. \\[1ex]

\hline%------------------------------------------------------------
\sumheading  If the project is successful, how will it make a difference? \\
\hline%------------------------------------------------------------
 It will introduce the ease of automation to domestic greenhouses, demystifying the struggles of home farming while encouraging self-sufficiency. By promoting food independence, it can reduce reliance on the national food supply chain, empowering individuals and communities to grow their own fresh produce sustainably.\\[1ex]

\hline%------------------------------------------------------------
\sumheading  What are the risks to the project being a success? Why is it expected to be successful? \\
\hline%------------------------------------------------------------
 The system may not be efficient enough at controlling the desired parameters. There may be delays due to time management issues. However, through thorough research, modular design and allowing buffer periods for unexpected delays, the project shouldn’t be able to fail. \\[1ex]

\hline%------------------------------------------------------------
\sumheading  What contributions have/will other students made/make? \\
\hline%------------------------------------------------------------
 No other students have/ will contribute to this design. \\[1ex]

\hline%------------------------------------------------------------
\sumheading  Which aspects of the project will carry on after completion and why? \\
\hline%------------------------------------------------------------
 Further research on the economic feasibility of the project may be conducted.  \\[1ex]

\hline%------------------------------------------------------------
\sumheading  What arrangements have been/will be made to expedite continuation? \\
\hline%------------------------------------------------------------
  Thorough documentation of the design process will allow for continuation or improvements to be made. Modular design should allow for ease of maintenance as well as allow for the upgrading of any components if necessary.  \\[1ex]

\hline%------------------------------------------------------------
\end{longtable}

